```latex
\documentclass{article}
\usepackage{amsmath,amssymb,amsthm}

\newtheorem{theorem}{Theorem}
\newtheorem{lemma}{Lemma}
\newtheorem{proposition}{Proposition}

\DeclareMathOperator{\Teich}{Teich}
\DeclareMathOperator{\Mod}{Mod}
\newcommand{\Sph}{\mathbb{S}^2}
\newcommand{\Dcal}{\mathcal{D}}

\begin{document}

\title{A Dehn twist with infinite orbit on the deformation space}
\author{}
\date{}
\maketitle

Let $f:\Sph^2\to \Sph^2$ be a quadratic rational map with critical points
$c_1,c_2$. Suppose that $c_1$ belongs to a cycle of length $4$, and that
the critical value $v=f(c_2)$ does not belong to that cycle. Put
\[
p_i=f^{\circ i}(c_1)\quad (i=1,2,3),
\]
so that
\[
A=\{c_1,p_1,p_2,p_3\},\qquad B=A\cup\{v\}.
\]
Thus the dynamics on the four-cycle is
\[
c_1\longmapsto p_1\longmapsto p_2\longmapsto p_3\longmapsto c_1.
\]
The two critical values of $f$ are
\[
V(f)=\{f(c_1),f(c_2)\}=\{p_1,v\}.
\]

We shall prove the following statement.

\begin{theorem}
Let $\gamma\subset (\Sph^2,B)$ be a simple closed curve whose two
complementary disks contain respectively
\[
\{c_1,v\}
\qquad\text{and}\qquad
\{p_1,p_2,p_3\}.
\]
Then the square of the Dehn twist about $\gamma$ preserves
$\Dcal_f$, and its action on the set of connected components of
$\Dcal_f$ has an infinite orbit.
\end{theorem}

The proof uses the following standard facts from mapping class group theory.

\begin{lemma}
\label{lem:trivial-twist}
Let $X$ be a sphere with a finite set of marked points. If $\alpha$ is a
simple closed curve in $X$ that bounds a disk containing at most one
marked point, then the Dehn twist $T_\alpha$ is isotopic to the identity
relative to the marked points. In particular, $T_\alpha$ and all its
powers act trivially on the corresponding Teichm\"uller space.
\end{lemma}

\begin{proof}
This is the usual Alexander trick: a disk with at most one marked point
has trivial mapping class group relative to its boundary and marked point.
\end{proof}

\begin{lemma}
\label{lem:preimage}
Let $\gamma$ be as in the theorem. Then $f^{-1}(\gamma)$ is a single
simple closed curve $\delta$, and $\delta$ bounds a disk containing
exactly one point of $A$, namely $p_3=f^{\circ 3}(c_1)$.
\end{lemma}

\begin{proof}
The curve $\gamma$ separates the two critical values: one side contains
$v$, and the other side contains $p_1=f(c_1)$. Since $f$ has degree $2$,
a simple closed curve separating the two critical values has connected
preimage; hence $f^{-1}(\gamma)=\delta$ is a single simple closed curve.

Let $D$ be the complementary disk of $\gamma$ containing $\{c_1,v\}$.
Then $D$ contains exactly one critical value, namely $v$. Therefore its
preimage $f^{-1}(D)$ is a disk. This disk contains the critical point
$c_2$, because $f(c_2)=v$, and it contains all preimages of $c_1$.
Among the points of $A$, the only point mapping to $c_1$ is $p_3$, since
the cycle has length $4$. Thus
\[
f^{-1}(D)\cap A=\{p_3\}.
\]
Hence the boundary $\delta=f^{-1}(\gamma)$ bounds a disk containing
exactly one point of $A$.
\end{proof}

\begin{proposition}
\label{prop:preserves}
The mapping class $T_\gamma^2$ preserves $\Dcal_f$.
\end{proposition}

\begin{proof}
Let $g=T_\gamma^2$.

First consider the forgetful map $i_*:\Teich(B)\to \Teich(A)$.
After forgetting the point $v$, the curve $\gamma$ bounds a disk
containing exactly the point $c_1\in A$. By Lemma~\ref{lem:trivial-twist},
the Dehn twist $T_\gamma$, and hence also $g=T_\gamma^2$, is isotopic
to the identity in $\Mod(\Sph^2,A)$. Therefore
\[
i_*\circ g=i_*.
\]

Now consider the pullback map $f^*:\Teich(B)\to \Teich(A)$. By
Lemma~\ref{lem:preimage}, the curve $\delta=f^{-1}(\gamma)$ bounds a disk
containing exactly one point of $A$, namely $p_3$. In an annular
neighborhood of $\gamma$, the covering $f$ is locally of the form
$z\mapsto z^2$. Hence a full Dehn twist downstairs lifts to a half-twist
upstairs, and therefore the square $g=T_\gamma^2$ lifts to a full Dehn
twist about $\delta$. By Lemma~\ref{lem:trivial-twist}, this twist is
trivial in $\Mod(\Sph^2,A)$. Thus
\[
f^*\circ g=f^*.
\]

Let $\tau\in\Dcal_f$. Then $f^*(\tau)=i_*(\tau)$, and
\[
f^*(g\tau)=f^*(\tau)=i_*(\tau)=i_*(g\tau).
\]
Hence $g\tau\in\Dcal_f$. Since $g$ is a homeomorphism of $\Teich(B)$,
this proves that $g$ preserves $\Dcal_f$.
\end{proof}

It remains to prove the infinite orbit statement.

\begin{proposition}
\label{prop:infinite}
The action of $g=T_\gamma^2$ on the set of connected components of
$\Dcal_f$ has an infinite orbit.
\end{proposition}

\begin{proof}
Let $\tau_0\in\Dcal_f$ be the point corresponding to the original complex
structure and the original map $f$. Let $D_0$ be the connected component
of $\Dcal_f$ containing $\tau_0$.

By Proposition~\ref{prop:preserves}, $g$ preserves $\Dcal_f$, so
$g^n(D_0)$ is a connected component of $\Dcal_f$ for every $n\in\mathbb Z$.

The pure mapping class group of $(\Sph^2,B)$ acts freely on
$\Teich(B)$. Since $\gamma$ separates the five marked points into two
subsets of size $2$ and $3$, the Dehn twist $T_\gamma$ is nontrivial and
has infinite order in $\Mod(\Sph^2,B)$. Hence $g=T_\gamma^2$ has infinite
order, and the points
\[
g^n(\tau_0)\qquad (n\in\mathbb Z)
\]
are pairwise distinct.

Because $i_*\circ g=i_*$, all these points lie in the same fiber
\[
F:=i_*^{-1}\bigl(i_*(\tau_0)\bigr)
\]
of the forgetful map $i_*:\Teich(B)\to \Teich(A)$.

The image of $\Dcal_f$ in the moduli space is the one-dimensional moduli
space of quadratic rational maps with a marked critical $4$-cycle and a
free critical value. In particular, the restriction
\[
i_*|_{D_0}:D_0\longrightarrow \Teich(A)
\]
is not constant. Since $D_0$ is a one-dimensional complex-analytic
subset of $\Teich(B)$ and $i_*$ is holomorphic, a nonconstant holomorphic
map from a one-dimensional complex space to a one-dimensional complex
manifold has discrete fibers. Thus $D_0\cap F$ is discrete. The same is
true for every translate $g^n(D_0)$, because
$i_*\circ g^n=i_*$.

Now suppose, toward a contradiction, that the orbit
$\{g^n(D_0):n\in\mathbb Z\}$ of components were finite. Then infinitely
many of the distinct points $g^n(\tau_0)$ would lie in a single component
$D_j=g^m(D_0)$. Since all these points lie in the same fiber $F$, the
intersection $D_j\cap F$ would be infinite, contradicting the discreteness
of the fibers of $i_*|_{D_j}$.

Therefore the orbit of the component $D_0$ under $g$ is infinite.
\end{proof}

\end{document}
```